\section{Related work}

\subsection{Visuomotor policy}
Visuomotor policy is a robot learning framework that maps raw sensory observations, typically high-dimensional visual inputs and low-dimensional robotic states, directly into low-level control actions. Early approaches, such as action chunking (ACT)~\citep{zhao2023learning}, utilized a conditional variational autoencoder with Transformers~\citep{vaswani2017attention} to learn fine-grained bimanual manipulation by predicting future action sequences. Diffusion policy~\citep{dp} introduced a diffusion generative paradigm by modeling the action distribution as a score-based gradient field~\citep{ddpm, song2020score}, significantly improving stability in multi-modal environments. Flow matching \citep{flowmatching} used in Vision-Language-Action (VLA) models, like $\pi$ family~\citep{intelligence2025pi,intelligence2025pi05, black2024pi_0}, has sought to simplify the generative process by learning straight-line probability paths between noise and action distributions. More recently, \citet{pan2025much} suggest that the selection of advanced model architectures (e.g., DiT~\citep{Peebles_2023_ICCV} and UNet\citep{dp}) and action chunking strategy~\citep{zhao2023learning} has a more significant impact on the success of flow matching than the regression method.

Despite these successes, diffusion methods suffer from high computational costs due to their iterative multi-step inference nature and complex architectures. Conversely, the vision-to-action model (VITA)~\citep{gao2025vita} pursues architectural minimalism by employing a lightweight Multi-Layer Perceptron (MLP) backbone for direct vision-to-action mapping. However, VITA relies completely on visual inputs, making it vulnerable to environmental visual distractors. Furthermore, VITA exhibits a dimensionality mismatch within its flow space. Recent efforts such as MeanFlow-based policies~\citep{fang2025omp, geng2025mean} and consistency models~\citep{zhang2025flowpolicy,li2026ofp, Consistency} have begun to push one-step generation into robotic action prediction by optimizing the theoretical boundaries of the noise-to-action mapping. In contrast, \m{} exploits the physical continuity of actions to directly map proprioceptive historical actions to future actions, fundamentally shortening the generative path and enabling high-performance single-step generation.

%While one-step inference has been extensively studied in image generation~\citep{lu2026one, Consistency, geng2025mean, kornilov2024optimal}, its adaptation to robotics remains limited.  
\subsection{Noise optimization of diffusion}

In image and video synthesis, optimizing initial noise is a key strategy for enhancing generation quality and speed~\citep{ahn2024noise,mao2024lottery,samuel2024generating,eyring2024reno, warmstarts}. For example, \citet{ahn2024noise} eliminate the need for computationally expensive guidance techniques by learning to map standard \textit{Gaussian} noise to a guidance-free noise space through a lightweight LoRA module~\citep{hu2022lora}. Similarly, warm-start diffusion~\citep{warmstarts} utilizes a deterministic model to provide an informed mean and variance of the initial \textit{Gaussian} distribution based on context, reducing the sampling path. In robotics, however, noise optimization is rarely explored. \citet{steeringdp} adapts behavioral cloning policies by running reinforcement learning over the latent-noise space, allowing the agent to steer robot actions without altering the base policy weights.

Departing from the long transport paths of the noise-to-data paradigm, \citet{chen2024bridger} leverage informative initial distributions, yet remains constrained by its dependence on heuristic or pre-trained models. Our \m{} introduces an action-to-action transport mechanism. By directly employing clean historical data as the initial distribution, we leverage the physical continuity of robotic motion to place the starting point much closer to the target in the high-dimensional latent space. This reduced distributional gap bypasses the need for multi-step refinement, enabling high-fidelity single-step inference optimized for real-time robotic control. 